\chapter{Monocular Depth Estimation}
\label{ch:mde}
\section{Introduction to Monocular Depth Estimation}
Depth estimation is the task of recovering, for every pixel in an image,
the distance between the camera and the corresponding point in the physical scene.
Traditionally, this information has been obtained through active sensing hardware,
such as LiDAR, ultrasonic rangefinders, or time-of-flight cameras, or through stereo vision,
where two calibrated cameras triangulate depth from the disparity between matched pixels. 
Depth estimation is a crucial component in many robots' navigation and [change wording] actuation because it allows them to 
perceive their environment and  


Monocular depth estimation (MDE), in contrast, seeks to recover this same per-pixel depth map
from a single RGB image and no additional hardware. In this sense, the goal of MDE is most similar
to that of an RGB-D camera, which uses an infrared projector to cast a structured grid of laser dots
onto the scene and infers per-pixel depth from how that pattern is distorted by scene geometry; MDE
aims to recover the same dense, per-pixel depth map, but passively, from RGB pixels alone and without
any active illumination. Because a single 2D projection is inherently
a many-to-one mapping of the 3D world, this problem is fundamentally ill-posed: an image of a small
object placed close to the camera is, without additional context, indistinguishable from an image
of a larger object placed further away. Resolving this ambiguity requires the model to exploit
monocular depth cues, e.g., relative object size, texture gradients, occlusion boundaries,
and learned priors about typical scene geometry, rather than direct geometric triangulation [\todo{cite?}].

Early work on monocular depth estimation used MRFbased formulations [3], simple geometric assumptions [2], or non-parametric methods [14] [midas paper]


Monocular depth estimation is particularly attractive for embedded and resource-constrained
robotic platforms, such as the cyclops rover, because it removes the need for a
dedicated depth sensor, stereo camera rig, or LiDAR unit, all of which add cost, weight,
and power draw that are difficult to justify on a small, battery-powered platform.
A single front-mounted camera can instead serve double duty: providing both the RGB imagery
used for scene understanding and, via MDE, an estimate of obstacle distance that would
otherwise require a separate range-finding sensor. Section~\ref{sec:mde-to-range} discusses
how this dense depth output is converted into the sparse range measurements consumed by the
localization pipeline described in Chapter~\ref{ch:localization}.

\subsection{Problem Formulation and Notation}
\label{sec:mde-formulation}
Let $I \in \mathbb{R}^{H \times W \times 3}$ denote an RGB image of height $H$ and width $W$,
and let $Z \in \mathbb{R}_{+}^{H \times W}$ denote the corresponding metric depth map, whose
entry $z(u,v)$ is the distance in meters from the camera's optical center to the scene point
imaged at pixel $(u,v)$. Monocular depth estimation seeks a function
\begin{equation}
  f_\theta : \mathbb{R}^{H \times W \times 3} \rightarrow \mathbb{R}^{H \times W},
  \qquad \hat{D} = f_\theta(I),
  \label{eq:mde-map}
\end{equation}
parameterized by learned weights $\theta$, that approximates $Z$ from $I$ alone. Throughout this
thesis $z$ denotes a true metric depth and $\hat{d}$ a raw network prediction; as
Section~\ref{sec:relative-depth} explains, these are not the same quantity and are not even
expressed in the same units.

The camera is modeled as an ideal pinhole with intrinsic matrix
\begin{equation}
  K = \begin{bmatrix}
        f_x & 0   & c_x \\
        0   & f_y & c_y \\
        0   & 0   & 1
      \end{bmatrix},
  \label{eq:intrinsics}
\end{equation}
where $f_x, f_y$ are the focal lengths in pixels and $(c_x, c_y)$ is the principal point. A scene
point $(X, Y, Z)$ expressed in the camera frame projects to
\begin{equation}
  u = f_x \frac{X}{z} + c_x,
  \qquad
  v = f_y \frac{Y}{z} + c_y,
  \label{eq:projection}
\end{equation}
and, given a depth $z$ at pixel $(u,v)$, the corresponding scene point is recovered by
back-projection:
\begin{equation}
  X = \frac{(u - c_x)\, z}{f_x},
  \qquad
  Y = \frac{(v - c_y)\, z}{f_y},
  \qquad
  Z = z .
  \label{eq:backprojection}
\end{equation}
Equation~\eqref{eq:backprojection} is the mechanism by which Section~\ref{sec:mde-to-range}
converts a depth map into the geometry of a range scan.

Equation~\eqref{eq:projection} also makes the ill-posedness described above precise. Consider
scaling the entire scene about the optical center by a factor $s > 0$, so that every scene point
$(X, Y, Z)$ is mapped to $(sX, sY, sZ)$. The projected coordinates become
\begin{equation}
  u' = f_x \frac{sX}{sz} + c_x = u,
  \qquad
  v' = f_y \frac{sY}{sz} + c_y = v,
  \label{eq:scale-ambiguity}
\end{equation}
which are unchanged, while every depth has been multiplied by $s$. The image is therefore
\emph{invariant} to a uniform rescaling of the scene, even though the depth map is not. No
function of $I$ alone, however expressive, can distinguish these cases; absolute scale is not
merely difficult to estimate from a single image, it is not present in the image at all.

\subsection{Relative versus Metric Depth}
\label{sec:relative-depth}
A consequence of Equation~\eqref{eq:scale-ambiguity} that is easy to overlook is that most
off-the-shelf MDE models do not, in fact, output meters. They output \emph{relative} depth: a
dense map that is unitless, typically normalized to $[0,1]$ over the image, and correct in
\emph{ordering} rather than in \emph{measurement}. A predicted value of $0.8$ asserts that the
pixel is nearer than one predicted at $0.4$; it does not assert that the pixel is $80$
centimeters away, and rescaling the whole prediction by any positive constant would be an equally
valid answer. For a robot, this is the wrong half of the problem to have solved: knowing which
obstacle is nearest is of little use to a motion planner that must decide whether there is enough
clearance to pass between two of them.

The first reason for this is geometric, and it is not a deficiency that better training could
remove. Figure~\ref{fig:scale-ambiguity} illustrates the argument of
Equation~\eqref{eq:scale-ambiguity} directly: an object of size $s$ at distance $z$ and an object
of size $2s$ at distance $2z$ subtend the same visual angle, fill the same cone of rays, and
therefore produce pixel-for-pixel identical images. Any model that maps images to depths must
assign both cases the same output, so at most one of the two answers can be correct in meters.
Recovering absolute scale requires information from outside the image, such as a known object
size, a known camera mounting height, or a second calibrated view.

The second reason is a consequence of how modern zero-shot models are trained. To generalize
across scenes rather than to a single dataset, models such as MiDaS, DPT, and Depth Anything are
trained on a mixture of heterogeneous sources: LiDAR-supervised driving data, RGB-D indoor
captures, stereo photographs scraped from the web, structure-from-motion reconstructions, and
frames extracted from 3D films \todo{cite MiDaS (Ranftl et al.), DPT, Depth Anything}. These
sources disagree about camera intrinsics, and several of them supply depth only up to an unknown
scale, or as raw stereo disparity with an unknown baseline, so depths drawn from different
datasets cannot be placed on a common axis at all. The standard remedy, introduced with MiDaS, is
a \emph{scale-and-shift-invariant} loss computed on inverse depth: before the error is measured,
the prediction is aligned to the ground truth by the scale and shift that best fit the two, and
only the residual is penalized. The practical effect is that the training signal never asks the
network to produce meters, because any prediction differing from the truth by an affine transform
of inverse depth is already scored as perfect. Working in inverse depth carries the additional
benefit of keeping sky and far-field pixels bounded, since $1/z \rightarrow 0$ rather than
diverging as $z$ grows.

What such a model does provide is an estimate related to true depth by an unknown affine
transform in the inverse-depth domain,
\begin{equation}
  \frac{1}{z(u,v)} \;\approx\; a\,\hat{d}(u,v) + b,
  \label{eq:affine-invariance}
\end{equation}
where the scale $a$ and shift $b$ are unknown and in general vary from image to image.
Equation~\eqref{eq:affine-invariance} is the useful statement of the limitation: the prediction is
not wrong, merely uncalibrated, and two correspondences of known metric depth suffice to determine
$a$ and $b$ and thereby convert the entire map to meters. Estimating them is the subject of
Section~\ref{sec:mde-to-range}.

\begin{figure}[htbp]
  \centering
  % Scale ambiguity in monocular depth estimation.
%
% An object of size s at distance z and an object of size 2s at distance 2z
% subtend exactly the same visual angle, and therefore produce exactly the
% same image. Used in section 3.1.2 to motivate relative depth output.
%
% Pure vector TikZ -- solid fills and plain paths only. Do NOT introduce
% opacity, shadings or fadings here: the Makefile builds via
% latex -> dvips -> dvipdf, where those features are unreliable.
%
% Geometry: the object-subtending rays have slope 0.18 and the field-of-view
% rays slope 0.35. The near box's front face sits at x = 3.2 (half-height
% 0.576) and the far box's at x = 6.4 (half-height 1.152) -- exactly twice
% the size at exactly twice the distance.
\begin{tikzpicture}[>=Latex, line join=round, font=\small]

  % ---- camera ---------------------------------------------------------
  \fill[black!70] (-0.45,-0.30) rectangle (0,0.30);
  \fill[black!70] (0,0) -- (0.35,0.45) -- (0.35,-0.45) -- cycle;
  \node[anchor=north] at (0,-0.55) {camera};

  % ---- field of view --------------------------------------------------
  \draw[black!45,dashed] (0,0) -- (7.6,2.66);
  \draw[black!45,dashed] (0,0) -- (7.6,-2.66);
  \node[black!55,anchor=south east] at (7.6,2.66) {field of view};

  % ---- rays subtending both objects -----------------------------------
  \draw[black!75] (0,0) -- (7.5,1.35);
  \draw[black!75] (0,0) -- (7.5,-1.35);

  % ---- hypothesis A: small object, near -------------------------------
  \fill[black!15]  (3.2,-0.576) rectangle (3.6,0.576);
  \draw[black!75]  (3.2,-0.576) rectangle (3.6,0.576);
  \node[anchor=south] at (3.4,0.70) {size $s$ at $z$};

  % ---- hypothesis B: twice the size, twice as far ----------------------
  \fill[black!28]  (6.4,-1.152) rectangle (7.2,1.152);
  \draw[black!75]  (6.4,-1.152) rectangle (7.2,1.152);
  \node[anchor=south] at (6.8,1.28) {size $2s$ at $2z$};

  \node[anchor=center] at (5.0,0) {\emph{or}};

  % ---- distance measures ----------------------------------------------
  \draw[black!35,dotted] (0,-0.55)   -- (0,-3.95);
  \draw[black!35,dotted] (3.2,-0.62) -- (3.2,-2.35);
  \draw[black!35,dotted] (6.4,-1.20) -- (6.4,-3.05);
  \draw[|<->|,black!65] (0,-2.25) -- (3.2,-2.25)
        node[midway,fill=white,inner sep=1.5pt] {$z$};
  \draw[|<->|,black!65] (0,-2.95) -- (6.4,-2.95)
        node[midway,fill=white,inner sep=1.5pt] {$2z$};

  % ---- the image -------------------------------------------------------
  \draw[->,black!60,line width=0.9pt] (8.15,0) -- (8.95,0);

  \draw[black!75,fill=black!3] (9.2,-1.6) rectangle (12.2,1.6);
  \fill[black!22] (10.1,-0.82) rectangle (11.3,0.82);
  \draw[black!75] (10.1,-0.82) rectangle (11.3,0.82);
  \node[anchor=north,align=center] at (10.7,-1.75)
        {the image $I$ produced\\by \emph{either} object};

\end{tikzpicture}

  \caption[Scale ambiguity in monocular depth estimation]{Scale ambiguity in monocular depth
    estimation. An object of size $s$ at distance $z$ and an object of size $2s$ at distance $2z$
    subtend the same visual angle and fill the same cone of rays, so both produce the identical
    image $I$. Since a uniform rescaling of the scene leaves the projection of
    Equation~\eqref{eq:projection} unchanged while multiplying every depth by the same factor,
    absolute scale is not recoverable from a single image. Standard monocular models therefore
    predict depth only up to an unknown scale and shift.}
  \label{fig:scale-ambiguity}
\end{figure}

\subsection{Evaluating Depth Predictions}
\label{sec:mde-metrics}
Depth predictions are conventionally scored against ground truth over the set of $N$ pixels for
which a valid ground-truth depth exists. Writing $\hat{z}$ for a predicted metric depth and $z$
for the corresponding ground truth, the metrics reported in Chapter~\ref{ch:results} are the
absolute relative error,
\begin{equation}
  \mathrm{AbsRel} = \frac{1}{N} \sum_{i=1}^{N} \frac{\lvert \hat{z}_i - z_i \rvert}{z_i},
  \label{eq:absrel}
\end{equation}
the root mean squared error in meters,
\begin{equation}
  \mathrm{RMSE} = \sqrt{\frac{1}{N} \sum_{i=1}^{N} (\hat{z}_i - z_i)^2},
  \label{eq:rmse}
\end{equation}
and the threshold accuracies $\delta_\tau$, defined as the fraction of pixels whose prediction
falls within a multiplicative tolerance of the truth,
\begin{equation}
  \delta_\tau = \frac{1}{N} \left\lvert \left\{ i : \max\!\left( \frac{\hat{z}_i}{z_i},
    \frac{z_i}{\hat{z}_i} \right) < \tau \right\} \right\rvert,
  \qquad \tau \in \{1.25,\ 1.25^2,\ 1.25^3\}.
  \label{eq:delta}
\end{equation}
AbsRel weights each error by its distance and so is dominated by the near field, whereas RMSE is
dominated by the far field; reporting both, together with $\delta_\tau$, gives a fuller picture
than any one of them alone.

Equations~\eqref{eq:absrel}--\eqref{eq:delta} all presume that $\hat{z}$ is expressed in meters,
and so cannot be applied to a relative-depth model as it stands. The standard practice is to align
each prediction to the ground truth before scoring it. For a scale-invariant model this means
median scaling,
\begin{equation}
  \hat{z} \leftarrow \hat{z} \cdot \frac{\mathrm{median}(z)}{\mathrm{median}(\hat{z})},
  \label{eq:median-scaling}
\end{equation}
while for an affine-invariant model the scale and shift of
Equation~\eqref{eq:affine-invariance} are fitted per image by least squares in the inverse-depth
domain. This alignment must be stated whenever such numbers are reported, because it grants the
model the very information, absolute scale, that the deployed system will have to obtain for
itself; a favorable AbsRel computed after median scaling says nothing about whether the rover can
measure a doorway.

For that reason, Chapter~\ref{ch:results} reports these full-frame metrics chiefly to situate the
model against the published literature, and treats them as secondary. The quantity that governs
whether localization succeeds is the error in the handful of range values that
Section~\ref{sec:mde-to-range} extracts from the depth map, measured after the scale recovery the
rover actually performs at run time.

\section{Vision Transformer Architecture in Monocular Depth Estimation}
\label{sec:vit-mde}

\subsection{From Convolutions to Global Attention}
A few popular architecture for MDE models are Vision Transformer-based model and CNN-based model. While both

The problem, why CNN don't work because of its inductive bias: locality. Where you sort of need to see everything. Comes: Vision Transformer. Vision Transformer was developed following the architecture of NLP models for transformer. The ViT quickly became the standard
\todo{Sharpen the locality argument with a concrete example: the apparent size of a doorway is
only interpretable relative to the room around it, so the cue that disambiguates scale may lie
far from the pixel being predicted, whereas a CNN must stack many layers before two distant
regions can interact. Ties back to the ``learned priors about typical scene geometry'' claim in
Section~\ref{sec:mde-formulation}.}

\subsection{Patch Embedding and Positional Encoding}
\todo{An image of size $H \times W$ is split into $N = HW/p^2$ non-overlapping patches of side
$p$, each flattened and linearly projected to a $D$-dimensional token; learned positional
encodings restore spatial order; the readout/class token and what DPT does with it.}

\subsection{Multi-Head Self-Attention}
\todo{Queries, keys and values; $\mathrm{softmax}(QK^\top/\sqrt{d_k})V$; multiple heads attending
to different relations. State the $O(N^2)$ cost in token count explicitly here, since it is what
ties input resolution to inference latency in Chapter~\ref{ch:embedded}.}

\subsection{Dense Prediction Transformer Decoder}
\todo{The piece that turns a classification backbone into a dense regressor: Reassemble
operations convert token sequences from several transformer stages back into spatial feature maps
at differing resolutions; Fusion blocks progressively upsample and merge them in a RefineNet-like
manner; a final head emits the dense inverse-depth map. Cite Ranftl et al., DPT.}

\subsection{The Model Deployed on the Cyclops Rover}
\todo{Name the specific checkpoint, input resolution, parameter count and quantization used on
the rover, and justify it against the embedded compute budget; forward-reference
Chapter~\ref{ch:embedded}.}

\section{Monocular Depth to Range Finder}
\label{sec:mde-to-range}

Relative depth is not sufficient for robot perception which requires exact metric distance to navigate. Some models do actually give that, DAV2 metric, though they are not at all accurate for high-precision robotics task. \todo{cite} et al, did this but not accurate. \todo{cite} et al did that but changes.

\subsection{Recovering Metric Scale}
\todo{Commit to one of: (i) known camera mounting height plus a flat-ground assumption, which
turns floor pixels into a metric constraint sufficient to solve for $a$ and $b$ in
Equation~\eqref{eq:affine-invariance}; (ii) a one-time calibration against a target at a known
distance; (iii) a metric model such as Depth Anything V2 metric, exploiting the fact that the
rover carries one fixed camera at one fixed mounting height, so the domain mismatch that
normally limits metric models does not apply. Then quantify the accuracy objection raised above
and give the failure modes of whichever method is chosen.}

\subsection{Back-Projection to 3D}
\todo{Apply Equation~\eqref{eq:backprojection} with the calibrated metric depth to lift each pixel
to a point $(X, Y, Z)$ in the camera frame, then transform into the rover body frame using the
fixed camera extrinsics.}

\subsection{Ground-Plane and Sky Rejection}
\todo{Floor pixels are close to the camera but are not obstacles; without rejecting them every
scan would report an imminent collision. Fit or assume a ground plane in the body frame and
discard points within a tolerance of it, along with points above the rover's height.}

\subsection{Reduction to a 1D Angular Range Scan}
\todo{The step that makes a camera behave like a range finder. The bearing of image column $u$ is
$\theta(u) = \arctan\!\big((u - c_x)/f_x\big)$; partition the horizontal field of view into $K$
angular bins matching the beam layout the measurement model of Chapter~\ref{ch:localization}
expects, and take the minimum surviving obstacle range within each bin. Define the maximum usable
range and the encoding of invalid or no-return beams.}

\subsection{Noise Characterization}
\todo{Characterize the error distribution of the resulting ranges as a function of true distance
and bearing, including the heavy tail produced by depth-boundary bleed at object silhouettes.
This distribution is the input to the beam model $p(z_t \mid x_t, m)$ in
Chapter~\ref{ch:localization}.}

%%% Local Variables: ***
%%% mode: latex ***
%%% TeX-master: "thesis.tex" ***
%%% End: ***
